CARA turns front-end uncertainty in rotation averaging from an input that is either trusted or ignored into a quantity that is recalibrated, tested against cycle closure, rescaled when the data drift, and converted into per-camera bounds with finite-sample guarantees. On synthetic scenes with overconfident, drifting uncertainty, it is more accurate than a strong robust baseline on all but sparse graphs, and its bounds stay close to their target under noise shifts. On 1DSfM Notre Dame, without front-end uncertainty, it gives bounds near their target and a shift signal, though not higher accuracy, and neither it nor any baseline survives structured doppelgangers.

With a learned front end on real images (VGGT on ETH3D), the bounds are close to their target on clean held-out scenes. Under rendered fog, low light and rain they fall short on hard cameras, and we trace this to camera-level failures that cancel in every cycle. Difficulty-conditional bounds raise coverage of the hardest cameras from 81.8\% to 92.8\% in a pre-registered run over 13 scenes, at the cost of flagging many correct cameras. A learned edge uncertainty separates outlier edges much better and raised the accuracy of both uncertainty-aware methods we tested, although not significantly after multiplicity correction; in an analysis that was not pre-registered, CARA's rotations and accepted edges raised the position accuracy after translation averaging. Accuracy gains over robust and confidence-weighted averaging remain small ($p=0.051$ against CG, uncorrected).

The next steps are a signal that detects misread images directly, larger view graphs, a second front end, and bundle adjustment under real adverse conditions.
